% Integrated FF.38--FF.39; review copy only, do not input separately.
\begin{claim}{FF.38}{Proved}\label{clm:FFnilpotentinputs}
The following two assertions hold.
\begin{enumerate}
\item Let \(\Gamma\) be a profinite group
nilpotent of some fixed finite class, and let
\(\chi:\Gamma\to\mathbb Q_p^\times\) be a
nontrivial continuous character. Then
\[
                 H^1(\Gamma,\mathbb Q_p(\chi))=0.
                                                        \tag{FF.38.1}
\]
\item Let \(S\) be a Hodge--Tate representation
of the Galois group of a finite extension
\(M/\mathbb Q_p\). If its inertia Lie algebra
is nilpotent, that Lie algebra is abelian.
After a finite extension of \(M\), the
inertia representation is semisimple and
has the canonical invariant projection
of Claim~\ref{clm:FFabelianHT}.
\end{enumerate}
\end{claim}
\begin{proof}
For the first assertion, induct on the
nilpotency class and put \(Z=Z(\Gamma)\), a closed central subgroup. If
\(\chi|_Z\ne1\), choose \(z\in Z\)
with \(\chi(z)\ne1\). For every cocycle
\(c\), centrality gives
\[
             (\chi(z)-1)c(g)=(\chi(g)-1)c(z).
\]
Thus \(c\) is a coboundary, proving
(FF.38.1). If \(\chi|_Z=1\), it descends
to a nontrivial character of \(\Gamma/Z\).
Inflation--restriction gives
\[
 0\longrightarrow H^1(\Gamma/Z,\mathbb Q_p(\chi))
 \longrightarrow H^1(\Gamma,\mathbb Q_p(\chi))
 \longrightarrow H^1(Z,\mathbb Q_p(\chi))^{\Gamma/Z}.
\]
The first group is zero by induction,
since \(\Gamma/Z\) has smaller
nilpotency class. The last group is
zero too: it is the invariant part
of \(\operatorname{Hom}_{\rm cont}(Z,\mathbb Q_p)
\otimes\mathbb Q_p(\chi)\), on which
\(\Gamma\) acts only through the
nontrivial scalar \(\chi\), because
\(Z\) is central. This proves the
first assertion, including the abelian
base case.

For the second, write \(\mathfrak g\)
for the inertia Lie algebra and
\(\Theta\in\mathfrak g\otimes\mathbb C_p\)
for the Sen operator. It is semisimple
by Hodge--Tate comparison
\citep[Section~2.7]{FF_Fontaine2004}.
Hence \(\operatorname{ad}\Theta\)
is semisimple on
\(\operatorname{End}(S)\otimes\mathbb C_p\),
and on its invariant subspace
\(\mathfrak g\otimes\mathbb C_p\).
It is nilpotent on that subspace
because \(\mathfrak g\) is nilpotent.
It must therefore be zero. Thus
\(\Theta\in Z(\mathfrak g)\otimes\mathbb C_p\).
The minimal-subspace assertion of
\citet[Theorem~3.1, \(d=0\)]{FF_Ohkubo2014}
implies
\(\mathfrak g\subset Z(\mathfrak g)\).
So \(\mathfrak g\) is abelian.

Matrix exponential identifies a
sufficiently small open subgroup
of the compact inertia image with
a lattice in this abelian Lie algebra.
That open subgroup is abelian.
A finite extension of \(M\) makes
its inertia image lie in this subgroup:
choose an open subgroup of \(G_M\)
whose intersection with inertia is
inside the given open inertia subgroup.
Claim~\ref{clm:FFabelianHT} applies
to the restricted Hodge--Tate
representation and proves the
remaining assertions. This concerns
a finite extension, not an infinite
passage that might destroy the
finite local-field hypotheses.
\end{proof}

\begin{claim}{FF.39}{Proved}\label{clm:FFnilpotentreconstruct}
Let \(I=G_{K_p}=I_{\mathbb Q_p}\),
and define its closed lower central
series by
\(\gamma_1(I)=I\) and
\(\gamma_{r+1}(I)=\overline{[I,\gamma_r(I)]}\).
For each integer \(c\ge1\), put
\[
              K_{p,c}=
                \overline{\mathbb Q}^{\,\gamma_{c+1}(I)}.
                                                        \tag{FF.39.1}
\]
Every finite extension of \(K_{p,c}\),
and every subfield of such an extension,
satisfies \'etale reconstruction.
Here \(K_{p,1}=K_p^{\rm ab}\), and
every extension \(K_{p,c+1}/K_{p,c}\)
is infinite. The nilpotency class
\(c\) is fixed and finite in the
assertion; no conclusion about the
union of these fields or a general
solvable closure is included.
\end{claim}
\begin{proof}
Write \(J_c=\gamma_{c+1}(I)\).
It is characteristic in inertia,
hence normal in \(G_{\mathbb Q_p}\).
The quotient \(I/J_c\) is nilpotent
of class at most \(c\).
The field in (FF.39.1) is exactly
the compositum of all finite Galois
extensions of \(K_p\) with Galois
group of nilpotency class at most
\(c\): the lower-central identities
pass to finite quotients, and the
intersection of the kernels of the
finite quotients of \(I/J_c\) is
trivial.

We first verify that these stages
are proper and form a strictly
increasing infinite hierarchy.
The free wild inertia group \(P\)
has \(H^2(P,\mathbb F_p)=0\).
The quotient \(I/P\) is pro-prime-to-\(p\),
whose positive-degree cohomology
on \(p\)-primary discrete modules
vanishes by averaging on finite
quotients. Hochschild--Serre gives
\[
                   H^2(I,\mathbb F_p)=0.
                                                        \tag{FF.39.2}
\]
Also \(H^1(I,\mathbb F_p)\) is
infinite-dimensional by the principal-unit
calculation in Step~1 of
Claim~\ref{clm:FFmaxabelianreconstruct}.

Every finite \(p\)-group \(E\)
is a quotient of \(I\). Here are
the embedding-problem details.
Its Frattini quotient \(E/\Phi(E)\)
is a finite-dimensional \(\mathbb F_p\)-space,
so the preceding infinite-dimensional
\(H^1\) supplies a surjection from
\(I\) to that quotient. Choose a
central chief series inside
\(\Phi(E)\): each successive
factor has order \(p\) and is
central in the corresponding quotient
of \(E\). Such a series exists
because a nontrivial normal subgroup
of a finite \(p\)-group meets
its center nontrivially, and one
continues in the quotient. Each
successive central extension by
\(\mathbb F_p\) has zero pullback
obstruction by (FF.39.2), so the
map from \(I\) lifts through all
of them. Its final image is all
of \(E\): a proper subgroup lies
in a maximal subgroup, and every
maximal subgroup contains
\(\Phi(E)\), contrary to surjectivity
on the Frattini quotient.

For each \(n\), apply this to
\(E=\operatorname{UT}_{c+2}(\mathbb F_p)^n\).
Each unitriangular factor has
nilpotency class \(c+1\), and
its \((c+1)\)-st lower-central
subgroup has order \(p\).
The nested commutator of the
successive elementary transvections
is \(1+E_{1,c+2}\), proving this
also in small characteristic.
It follows that
\(\gamma_{c+1}(I)/\gamma_{c+2}(I)\)
has quotients \((\mathbb Z/p)^n\)
for arbitrarily large \(n\).
It is infinite, proving the
asserted infinite stage extensions.

We will also need that \(J_c\)
and all its open subgroups are
infinite. Independently of the
preceding quotient argument,
\(J_c\) contains \(\gamma_{c+1}(P)\).
Map \(c+1\) free generators of
\(P\) to the successive elementary
transvections just used. Their
nested commutator is nonidentity,
so \(\gamma_{c+1}(P)\ne1\).
Torsion-freeness of \(P\) makes
it infinite. Every open subgroup
of \(J_c\) is infinite as well.

Now fix a finite extension
\(K/K_{p,c}\). It contains
\(K_p^{\rm ab}\), hence
\(U_p\) and \(\mu_\infty\).
Write \(K=L_0K_{p,c}\), with
\(L_0/K_p\) finite, enlarging
\(L_0\) to contain any specified
finite data. The Galois group
\(\operatorname{Gal}(K/L_0)\)
is a closed subgroup of the
nilpotent group \(I/J_c\),
so is nilpotent of class at most
\(c\), with infinite cyclotomic
image. Claim~\ref{clm:FFnilpotentinputs}(1)
replaces the abelian cocycle identity
in the proof of (FF.37.2).
It proves the same exact integral
and rational Kummer kernels
\(\mu_\infty\), because an
infinitely \(p\)-divisible element
has its compatible Kummer cocycle
factoring through this nilpotent
quotient. Rational triviality there
implies that a \(p\)-power of
the element belongs to the known
integral kernel over \(L_0\),
and hence the element is a
root of unity.

Choose a finite number model
\(D_0\subset K\), put
\(M_0=(D_0)_v\), and identify
\[
        G_K=H_0=J_c\cap G_{M_0}
                         \triangleleft G_{M_0}.
                                                        \tag{FF.39.3}
\]
The quotient \(I_{M_0}/H_0\)
is nilpotent of class at most
\(c\) and has nontrivial
cyclotomic character. Assertion~(1)
again gives zero \(H^1\) of this
quotient with coefficients
\(\mathbb Q_p(1)\), while its
coefficient invariants are zero.
The two inflation--restriction
arguments in (FF.37.4) prove
injectivity from both inertia
and the full local Galois group
to \(H_0\). The same statements
hold at every finite field and
finite local-model extension.

For a pointed generalized Albanese,
form the finite-dimensional space
\(S=\operatorname{Hom}_{H_0}(V_pA,\mathbb Q_p(1))\).
It is Hodge--Tate and its inertia
image is nilpotent of class at
most \(c\); its Lie algebra
is therefore nilpotent.
By Claim~\ref{clm:FFnilpotentinputs}(2),
a finite local extension makes
this inertia representation abelian
and semisimple. Realize that
extension on a finite algebraic
model, possibly passing to a
finite extension of \(K\).
Retain the original space \(S\)
under this restriction; it is
unnecessary to replace it by
a larger Hom space. The finite
extension is permissible in the
point-function descent argument.

The all-point-invisible subspace
\(N\subset S\) of Step~3 in
Claim~\ref{clm:FFmaxabelianreconstruct}
has the same definition and the
same normality and generation
properties. It is unchanged by
testing only finite extensions
of the enlarged base: any finite
extension of the original base
can be combined with that base,
and rational restriction--corestriction
preserves a nonzero evaluation.
The inertia detector of
Claim~\ref{clm:FFpdetector}, followed
by the just-proved restriction
injection, again gives
\(N\cap S^{I_{M_1}}=0\) at
every further finite local stage.
Actual point/value descent and
the canonical semisimple projection
therefore give the unique invariant
lift \(\lambda_0\) modulo \(N\),
exactly as in (FF.37.7)--(FF.37.8).
The projection remains compatible
with torus restrictions, open
restrictions and curve pullback;
in particular it retains the
original boundary residues.

All remaining inputs of the
reconstruction proof in Step~5
of Claim~\ref{clm:FFmaxabelianreconstruct}
have now been established: the
exact Kummer kernel, finite-local
restriction injection, unchanged
integral boundary residues, and
the infinite groups of every
finite extension needed for
large-degree points. Thus the
normalized rational class yields
a power of an algebraic unit,
the \(T,T-1\) resultant removes
that power, and \(T-\beta,T-a\)
remove the constant and root-of-unity
ambiguities in the actual point
function. Finite descent returns
the resulting unit to the
original \(K\), and
\citep[Theorem~4.18(5)]{FF_CHW2024}
gives \(\mathcal R(K)\).
Algebraic descent proves the
subfield assertion, since all
these fields are algebraic
over \(\mathbb Q\).
The proof uses a fixed bound
\(c\) in the cohomological
central-series induction and
the inertia Lie argument;
neither argument has been
applied to an unbounded union.
\end{proof}
